\section{Métricas}

La medición de la arquitectura de seguridad permite verificar que los controles y capacidades
definidos en las secciones anteriores produzcan los resultados esperados sobre la operación del CCI. Para ello se retoma el perfil de atributos de negocio
desarrollado mediante SABSA, en el cual cada requerimiento de seguridad se encuentra asociado
con indicadores de desempeño (\textit{Key Performance Indicators}, KPI) e indicadores de
riesgo (\textit{Key Risk Indicators}, KRI).

Los KPI representan las metas de desempeño que debe alcanzar la arquitectura y permiten
evaluar la efectividad de los controles implementados. Los KRI funcionan como mecanismos de
alerta temprana y establecen umbrales amarillo y rojo que facilitan identificar desviaciones antes
de que estas generen un nivel de exposición no aceptable.

Esta estructura permite relacionar directamente las métricas con los riesgos R1 a R4, así como
con los dominios de seguridad, servicios O-ESA y capacidades arquitectónicas implementadas
durante las diferentes fases de la hoja de ruta.


\subsection{Indicadores y Métricas de Desempeño (KPI / KRI)}

La Tabla~\ref{tab:kpi_kri_1} y la Tabla~\ref{tab:kpi_kri_2} consolidan los indicadores
definidos para evaluar la integridad y oportunidad de la telemetría, el control de acceso,
la trazabilidad de las decisiones, la disponibilidad y recuperación del servicio, el monitoreo
de seguridad, la protección de datos personales y el cumplimiento normativo.


\begin{table}[H]
	\centering
	\caption{Indicadores KPI y KRI de la arquitectura de seguridad del CCI (BA1--BA6)}
	\label{tab:kpi_kri_1}
	
	\begingroup
	\scriptsize
	\setstretch{1}
	\setlength{\tabcolsep}{3pt}
	\renewcommand{\arraystretch}{1.12}
	
	\begin{tabularx}{\textwidth}{
			>{\centering\arraybackslash}p{0.8cm}
			>{\raggedright\arraybackslash}p{2.0cm}
			>{\centering\arraybackslash}p{1.4cm}
			Y
			Y}
		
		\toprule
		\textbf{ID} &
		\textbf{Atributo} &
		\textbf{Riesgo} &
		\textbf{KPI (meta)} &
		\textbf{KRI (amarillo / rojo)} \\
		\midrule
		
		BA1 &
		Integridad asegurada &
		R1 &
		Al menos 99,5\% de los mensajes GPS/AVL con autenticidad e integridad verificadas. &
		Registros rechazados o anómalos por cada 100.000: amarillo $>$ 50; rojo $>$ 200. \\
		
		\addlinespace[2pt]
		
		BA2 &
		Exacta y actual &
		R1 &
		Latencia de posición p95 $\leq$ 30 s. &
		Buses sin reporte vigente por más de 2 min: amarillo $>$ 2\%; rojo $>$ 5\%. \\
		
		\addlinespace[2pt]
		
		BA3 &
		Autenticada y autorizada &
		R3 &
		100\% de cuentas de consolas críticas con MFA y 100\% de revisiones trimestrales
		de permisos ejecutadas. &
		Accesos fuera de perfil u horario por 1.000 sesiones: amarillo $>$ 1; rojo $>$ 5. \\
		
		\addlinespace[2pt]
		
		BA4 &
		Segregación de funciones &
		R2, R3 &
		Cero usuarios con roles incompatibles. &
		Excepciones de segregación abiertas: amarillo $\geq$ 1; rojo $\geq$ 3. \\
		
		\addlinespace[2pt]
		
		BA5 &
		Trazable y responsable &
		R2 &
		Al menos 99,9\% de decisiones con registro completo. &
		Decisiones sin registro completo por 1.000: amarillo $>$ 1; rojo $>$ 5. \\
		
		\addlinespace[2pt]
		
		BA6 &
		Auditable y no repudiable &
		R2 &
		100\% de verificaciones periódicas de integridad del repositorio WORM aprobadas. &
		Intentos de modificar o borrar registros: amarillo = 1 bloqueado; rojo = cualquier
		alteración efectiva. \\
		
		\bottomrule
	\end{tabularx}
	
	\endgroup
\end{table}


\begin{table}[H]
	\centering
	\caption{Indicadores KPI y KRI de la arquitectura de seguridad del CCI (BA7--BA11)}
	\label{tab:kpi_kri_2}
	
	\begingroup
	\scriptsize
	\setstretch{1}
	\setlength{\tabcolsep}{3pt}
	\renewcommand{\arraystretch}{1.12}
	
	\begin{tabularx}{\textwidth}{
			>{\centering\arraybackslash}p{0.8cm}
			>{\raggedright\arraybackslash}p{2.0cm}
			>{\centering\arraybackslash}p{1.4cm}
			Y
			Y}
		
		\toprule
		\textbf{ID} &
		\textbf{Atributo} &
		\textbf{Riesgo} &
		\textbf{KPI (meta)} &
		\textbf{KRI (amarillo / rojo)} \\
		\midrule
		
		BA7 &
		Disponible &
		R4 &
		Disponibilidad mensual $\geq$ 99,9\%. &
		Minutos de indisponibilidad en el mes: amarillo $>$ 20; rojo $>$ 43. \\
		
		\addlinespace[2pt]
		
		BA8 &
		Continua y recuperable &
		R4 &
		RTO $\leq$ 15 min y 100\% de pruebas semestrales de conmutación exitosas. &
		Días desde la última prueba exitosa: amarillo $>$ 150; rojo $>$ 180. \\
		
		\addlinespace[2pt]
		
		BA9 &
		Detectable y monitoreada &
		R1--R4 &
		Tiempo medio de detección $\leq$ 5 min en eventos críticos y al menos 95\% de
		activos críticos reportando al SIEM. &
		Alertas críticas sin atender por más de 15 min: amarillo $\geq$ 1; rojo $\geq$ 5. \\
		
		\addlinespace[2pt]
		
		BA10 &
		Confidencial y privada &
		Transversal &
		100\% de indicadores y tableros sin datos personales identificables innecesarios. &
		Accesos a datos personales sin justificación: amarillo $\geq$ 1; rojo $\geq$ 3. \\
		
		\addlinespace[2pt]
		
		BA11 &
		Cumplimiento &
		Transversal &
		Al menos 90\% de controles aplicables con evidencia vigente. &
		Hallazgos de auditoría vencidos: amarillo $\geq$ 1; rojo $\geq$ 3. \\
		
		\bottomrule
	\end{tabularx}
	
	\endgroup
\end{table}


Los indicadores permiten evaluar diferentes dimensiones de la arquitectura. BA1 y BA2
verifican la confiabilidad y oportunidad de la telemetría utilizada por el Sistema de Despacho;
BA3 y BA4 evalúan la efectividad del gobierno de identidades y la segregación de funciones;
BA5 y BA6 miden la trazabilidad, integridad y capacidad probatoria de la bitácora; mientras que
BA7 y BA8 permiten comprobar la disponibilidad y capacidad de recuperación del CCI.

De manera transversal, BA9 evalúa la capacidad de detección y monitoreo de eventos sobre los
activos críticos, mientras que BA10 y BA11 permiten verificar respectivamente la protección de
datos personales y el cumplimiento de los controles normativos aplicables.

El seguimiento de estos indicadores debe integrarse con los mecanismos de monitoreo y
auditoría definidos en la arquitectura objetivo. Los KRI que alcancen el nivel amarillo deben
generar una revisión preventiva del control asociado, mientras que un umbral rojo requiere
escalamiento y tratamiento prioritario por parte de los responsables de seguridad y operación.

Finalmente, las métricas deben utilizarse como insumo para revisar periódicamente el riesgo
residual de R1 a R4. De esta forma, la arquitectura no se considera efectiva únicamente por la
implementación de los mecanismos técnicos definidos en la hoja de ruta, sino por su capacidad
demostrable para mantener la integridad de los datos, limitar accesos no autorizados, preservar
la trazabilidad y sostener la disponibilidad de los servicios críticos del CCI.